\documentclass[a4paper, serif, 10pt]{moderncv}

%% moderncv
% style and color
\moderncvstyle{banking}
\moderncvcolor{black}

% renew moderncv command to adapt for CJK colon
\renewcommand*{\cvitem}[3][.25em]{%
  \ifstrempty{#2}{}{\hintstyle{#2}：}{#3}%
  \par\addvspace{#1}}

\renewcommand*{\cvitemwithcomment}[4][.25em]{%
  \savebox{\cvitemwithcommentmainbox}{\ifstrempty{#2}{}{\hintstyle{#2}：}#3}%
  \setlength{\cvitemwithcommentmainlength}{\widthof{\usebox{\cvitemwithcommentmainbox}}}%
  \setlength{\cvitemwithcommentcommentlength}{\maincolumnwidth-\separatorcolumnwidth-\cvitemwithcommentmainlength}%
  \begin{minipage}[t]{\cvitemwithcommentmainlength}\usebox{\cvitemwithcommentmainbox}\end{minipage}%
  \hfill% fill of \separatorcolumnwidth
  \begin{minipage}[t]{\cvitemwithcommentcommentlength}\raggedleft\small\itshape#4\end{minipage}%
  \par\addvspace{#1}}

%% page layout/margins
\usepackage[top=2.5cm, bottom=2.5cm, left=1.5cm, right=1.5cm]{geometry}

\nopagenumbers{}

%% Babel config for English language
\usepackage[english]{babel}

%% fontspec
\usepackage{fontspec}

\IfFontExistsTF{Linux Libertine}{
  \setmainfont[Ligatures={TeX, Common}, Numbers=Lining, ItalicFont=Linux Libertine]{Linux Libertine}
}{}
\IfFontExistsTF{Linux Libertine O}{
  \setmainfont[Ligatures={TeX, Common}, Numbers=Lining, ItalicFont=Linux Libertine O]{Linux Libertine O}
}{}

%% CTeX
% CJK support, used to show CJK characters in the resume
%
% - fontset=none: disable builtin fontset but instead set the CJK font manually
% - heading=false: disable ctex heading
% - punct=kaiming: use kaiming punctuations styles for CJK
% - scheme=plain: use plain scheme, do not override \normalsize font size
% - space=auto: space settings for CJK characters
%
% ref:
% - http://ctan.mirrorcatalogs.com/language/chinese/ctex/ctex.pdf
\usepackage[UTF8, heading=false, punct=kaiming, scheme=plain, space=auto]{ctex}

\IfFontExistsTF{Noto Serif CJK SC}{
  \setCJKmainfont{Noto Serif CJK SC}
}{}
\IfFontExistsTF{Noto Sans CJK SC}{
  \setCJKsansfont{Noto Sans CJK SC}
}{}

%% line spacing
\usepackage{setspace}
\setstretch{1.125}

%% URL styling - use normal text instead of monospace
\urlstyle{same}

% Auto-underline all links
% Defer until \begin{document} so hyperref is already loaded
\AtBeginDocument{%
  \let\oldhref\href
  \renewcommand{\href}[2]{\oldhref{#1}{\underline{#2}}}%
}

%% Patch \httplink and \httpslink to handle full URLs
% The original moderncv macros always prepend http:// or https:// to URLs.
% This causes issues when the URL already has a protocol (e.g., https://example.com)
% resulting in malformed URLs like https://https://example.com.
% This patch checks if the URL already contains "://" and uses it directly if so.
% Note: We use \str_set:Nx (x-expansion) to expand macros like \@homepage first.
\ExplSyntaxOn
\str_new:N \g_yamlresume_url_str

\RenewDocumentCommand{\httpslink}{O{}m}{
  \str_set:Nx \g_yamlresume_url_str {#2}
  \str_if_in:NnTF \g_yamlresume_url_str {://}
    { \url{#2} }
    { \url{https://#2} }
}

\RenewDocumentCommand{\httplink}{O{}m}{
  \str_set:Nx \g_yamlresume_url_str {#2}
  \str_if_in:NnTF \g_yamlresume_url_str {://}
    { \url{#2} }
    { \url{http://#2} }
}
\ExplSyntaxOff

\name{陳嘉俊}{}
\title{資深 DevOps 工程師 | 雲端架構與自動化交付專家}
\phone[mobile]{+852 9123 4567}
\email{kachun.chan@example.com}
\homepage{https://kachunchan.dev}

\address{中國香港，九龍，香港，香港九龍觀塘鴻圖道 88 號創紀大廈 12 樓}{}{}

\extrainfo{{\small \faGithub}\ \href{https://github.com/kachunchan}{@kachunchan} {} {} {} • {} {} {} 
{\small \faLinkedin}\ \href{https://www.linkedin.com/in/kachunchan}{@kachunchan} {} {} {} • {} {} {} 
{\small \faStackOverflow}\ \href{https://stackoverflow.com/users/1842371/kachunchan}{@kachunchan}}

\begin{document}

\maketitle

\section{簡介}

\cvline{}{\begin{itemize}
\item 8 年雲端基礎架構與 DevOps 經驗，專注於大規模容器化平台、CI/CD 自動化與可觀測性建設
\item 主導單體式系統重構為 Kubernetes 微服務架構，部署頻率由每月 2 次提升至每日 10 次以上
\item 以 Terraform 與 Ansible 全面實踐基礎架構即程式碼（IaC），環境建置時間由 3 個工作日縮短至 45 分鐘
\item 建立 Prometheus、Grafana 與 ELK 監控堆疊，平均故障恢復時間（MTTR）下降 70\%，系統可用率達 99.98\%
\item 雲端成本治理經驗豐富，透過資源調度與預留實例策略，每年節省逾 180 萬港元雲端支出
\end{itemize}
熱衷推動 DevOps 文化與 blameless 事後檢討機制，擅長跨部門溝通協作，並定期主持內部技術工作坊與新人培訓。}
\section{教育背景}

\cventry{2012年9月–2016年6月}
        {學士，計算機科學，成績：3.6}
        {香港科技大學}
        {\url{https://www.ust.hk}}
        {}
        {\begin{itemize}
\item 主修計算機科學，成績位列全系前 15\%，連續三年獲院長嘉許名單
\item 畢業專題為以容器技術建置校園課程排程系統，獲得學系最佳專題獎
\item 擔任計算機學會幹事，籌辦校內黑客松與技術講座
\end{itemize}
\textbf{課程}：資料結構與演算法、作業系統、電腦網絡、資料庫系統、分散式系統、雲端運算導論、資訊安全基礎}
\section{工作經歷}

\cventry{2021年4月至今}
        {資深 DevOps 工程師}
        {雲際科技（香港）有限公司}
        {\url{https://www.cloudhorizon.com.hk}}
        {}
        {\begin{itemize}
\item 帶領 5 人平台工程小組，負責 12 條產品線的 Kubernetes（Amazon EKS）多集群架構設計、升級與維運
\item 以 Terraform 模組化管理 AWS 資源（VPC、EKS、RDS、IAM、S3），將新環境交付時間由 3 個工作日縮短至 45 分鐘
\item 建置 GitLab CI 與 Argo CD 的 GitOps 流水線，實現每日 10 次以上零停機部署，回滾時間少於 3 分鐘
\item 導入 Prometheus、Grafana、Loki 與 OpenTelemetry，建立 SLO 告警與錯誤預算機制，MTTR 由 90 分鐘降至 25 分鐘
\item 推動雲端成本治理（FinOps），透過 Karpenter 自動擴縮與 Savings Plans 規劃，年度雲端支出下降 32\%
\item 主持每月內部技術工作坊及事故事後檢討，推動知識文件化與值班輪替制度
\end{itemize}
\textbf{關鍵字}：Kubernetes、Terraform、GitOps、Observability、FinOps}

\cventry{2018年7月–2021年3月}
        {DevOps 工程師}
        {亞太數據服務有限公司}
        {\url{https://www.apacdata.com.hk}}
        {}
        {\begin{itemize}
\item 將單體式 Java 應用拆分為 20 個微服務並容器化部署至 Kubernetes，發佈週期由季度縮短至雙週
\item 以 Jenkins 與 Ansible 建立自動化建置、測試與部署流程，人工操作步驟減少 85\%
\item 設計香港與新加坡多區域災難復原架構，RTO 由 8 小時降至 30 分鐘，RPO 控制在 5 分鐘內
\item 導入 HashiCorp Vault 集中管理密鑰與憑證，消除設定檔中的明文密碼與共用帳號
\item 與開發團隊協作推行測試自動化與品質閘門，單元測試覆蓋率由 42\% 提升至 81\%
\end{itemize}
\textbf{關鍵字}：Jenkins、Ansible、Docker、Vault、Disaster Recovery}

\cventry{2016年8月–2018年6月}
        {系統運維工程師}
        {凱達資訊科技顧問有限公司}
        {\url{https://www.kaidata-its.com}}
        {}
        {\begin{itemize}
\item 負責 200 多台 Linux 伺服器的日常維運、效能調校、修補管理與容量規劃
\item 建置 Zabbix 與 Nagios 監控告警系統，將被動故障處理轉為主動預警
\item 撰寫 Shell 與 Python 自動化腳本，將例行備份與巡檢工作壓縮至每日 15 分鐘內完成
\item 協助三家客戶完成 AWS EC2 與 RDS 遷移，平均節省 40\% 自建機房維運成本
\end{itemize}
\textbf{關鍵字}：Linux、Shell、Python、Zabbix}
\section{語言}

\cvline{粵語}{母語或雙語水平 \hfill \textbf{關鍵字}：母語}
\cvline{英語}{完全專業水平 \hfill \textbf{關鍵字}：IELTS 7.5、商務書信與技術簡報}
\cvline{普通話}{完全專業水平 \hfill \textbf{關鍵字}：普通話水平測試二級甲等}
\section{專業技能}

\cvline{雲端平台}{專家 \hfill \textbf{關鍵字}：AWS、Google Cloud、Azure、Amazon EKS、Lambda}
\cvline{容器與編排}{專家 \hfill \textbf{關鍵字}：Docker、Kubernetes、Helm、Istio、Karpenter}
\cvline{基礎架構即程式碼}{高級 \hfill \textbf{關鍵字}：Terraform、Ansible、Pulumi、CloudFormation}
\cvline{CI/CD 與 GitOps}{高級 \hfill \textbf{關鍵字}：GitLab CI、Jenkins、GitHub Actions、Argo CD}
\cvline{監控與可觀測性}{高級 \hfill \textbf{關鍵字}：Prometheus、Grafana、Loki、OpenTelemetry、ELK Stack}
\cvline{程式語言與腳本}{高級 \hfill \textbf{關鍵字}：Python、Go、Bash、SQL}
\section{榮譽}

\cventry{2023年12月}
        {雲際科技（香港）有限公司}
        {年度卓越工程團隊獎}
        {}
        {}
        {表揚在雲端成本治理專案中的傑出表現，透過自動擴縮與資源調度策略，於 12 個月內為公司節省逾 180 萬港元雲端支出，同時維持服務可用率 99.98\%。}

\cventry{2020年1月}
        {亞太數據服務有限公司}
        {最佳自動化專案獎}
        {}
        {}
        {肯定在 CI/CD 自動化改造專案中的貢獻，將部署流程由人工操作轉為全自動流水線，發佈失誤率下降 76\%。}
\section{證書}

\cventry{2022年3月}
        {Cloud Native Computing Foundation}
        {Certified Kubernetes Administrator (CKA)}
        {\url{https://www.cncf.io/certification/cka/}}
        {}
        {}

\cventry{2023年8月}
        {Cloud Native Computing Foundation}
        {Certified Kubernetes Security Specialist (CKS)}
        {\url{https://www.cncf.io/certification/cks/}}
        {}
        {}

\cventry{2021年11月}
        {Amazon Web Services}
        {AWS Certified DevOps Engineer – Professional}
        {\url{https://aws.amazon.com/certification/}}
        {}
        {}

\cventry{2023年5月}
        {HashiCorp}
        {HashiCorp Certified: Terraform Associate}
        {\url{https://www.hashicorp.com/certification}}
        {}
        {}
\section{出版刊物}

\cventry{2023年8月}
        {以 GitOps 實現多集群 Kubernetes 的零停機交付}
        {香港雲端原生技術年刊 2023}
        {\url{https://www.cloudnativehk.org/journal/2023/gitops-multi-cluster}}
        {}
        {\begin{itemize}
\item 探討以 Argo CD 與 ApplicationSet 管理跨區域多集群部署的實務模式
\item 比較 push-based 與 pull-based 交付流程在合規與安全層面的差異
\item 分享將回滾時間由 20 分鐘壓縮至 3 分鐘的具體調校經驗
\end{itemize}}
\section{引薦}

\cventry{\emaillink[wendy.lee@example.com]{wendy.lee@example.com}}
        {雲際科技（香港）有限公司 平台工程總監}
        {李婉婷}
        {+852 2765 4321}
        {}
        {嘉俊在任職期間主導了我們整個 Kubernetes 平台的現代化工程，他不僅技術紮實，更具備難得的推動力與溝通能力，能將複雜的架構議題轉化為團隊可執行的路線圖，是任何工程組織都值得擁有的人才。}
\section{項目}

\cventry{2022年3月至今}
        {開源的 Kubernetes 資源使用率分析與成本優化工具，協助團隊找出閒置工作負載並提出調整建議}
        {KubeSaver}
        {\url{https://github.com/kachunchan/kubesaver}}
        {}
        {\begin{itemize}
\item 以 Go 語言開發，定期抓取 Prometheus 指標並計算每個命名空間的成本分攤
\item 提供 CLI 與網頁儀表板，可輸出建議的 requests/limits 調整清單
\item 累積超過 1,200 個 GitHub Star，並獲多間香港初創企業採用於內部成本治理
\end{itemize}
\textbf{關鍵字}：Kubernetes、FinOps、Go、Prometheus}

\cventry{2021年9月–2022年6月}
        {為公司內部 30 多個專案提供一致化的 CI/CD 樣板與安全掃描閘門}
        {CICD 樣板庫}
        {\url{https://github.com/kachunchan/pipeline-templates}}
        {}
        {\begin{itemize}
\item 整合 Trivy、Semgrep 與 SBOM 產生流程，於流水線中自動攔截高風險漏洞
\item 將新專案導入 CI/CD 的時間由 2 日縮短至 30 分鐘內完成
\end{itemize}
\textbf{關鍵字}：GitLab CI、DevSecOps、Trivy}
\section{興趣愛好}

\cvline{運動}{長跑、龍舟、游泳}
\cvline{個人興趣}{攝影、開源社群、手沖咖啡}
\section{志願服務}

\cventry{2021年3月至今}
        {講者及工作坊導師}
        {香港雲端原生社群}
        {\url{https://www.cloudnativehk.org}}
        {}
        {\begin{itemize}
\item 每季主持 Kubernetes、CI/CD 及可觀測性主題工作坊，累計參與人數超過 600 人
\item 協助籌辦年度香港 DevOps 年會，負責議程規劃、講者邀請與現場技術支援
\item 為初創團隊提供雲端架構諮詢，協助其在有限預算下建立可靠的自動化交付流程
\end{itemize}}
    
\end{document}